\documentclass[../../main.tex]{subfiles}
\begin{document}
\begin{oframed}
  \subfile{problem}
\end{oframed}

{\bf [解答]}

$\alpha = a+b, \beta = ab$ とすると，与えられた不等式を変形して
\begin{align}
  & 3 \left( \dfrac{\alpha+c}{3}-\sqrt[3]{c\beta} \right) - 2 \left( \dfrac{\alpha}{2}-\sqrt{\beta} \right) \ge 0 \\
  \therefore
  (\alpha+c) - 3\sqrt[3]{\beta} \cdot c^{\frac{1}{3}} - (\alpha - 2\sqrt{\beta}) \ge 0 \label{eq:1}
\end{align}
である．この左辺を$f(c)$とおく：
\begin{align}
  f(c) =  (\alpha+c) - 3\sqrt[3]{\beta} \cdot c^{\frac{1}{3}} - (\alpha - 2\sqrt{\beta})
\end{align}
とおくと，一階微分は
\begin{align}
  f'(c) = 1 - \sqrt[3]{\beta} \cdot c^{-\frac{2}{3}}
\end{align}
だから，$f(c)$の増減表は以下のようになる．

\begin{table}[htb]
  \centering
  \caption{$f(c)$の増減表}
  \begin{tabular}{|c|c|c|c|c|}
    \hline
    $c$ & $0$ & $\cdots$ & $\sqrt{|\beta|}$ & $\cdots$ \\
    \hline
    $f'$ & & $-$ & $0$ & $+$ \\
    \hline
    $f$ & & $\searrow$ & & $\nearrow$ \\
    \hline
  \end{tabular}
\end{table}

よって,
\begin{align}
  f(c) \ge f(\sqrt{|\beta|}) = \alpha + \sqrt{|\beta|} - 3\sqrt{|\beta|} - \alpha + 2\sqrt{|\beta|} = 0
\end{align}
となり\cref{eq:1}は示された．

等号成立は$c=\sqrt{\beta}=\sqrt{ab}$の時である．

\end{document}
